% =======================================================================
% Code-Listings Konfiguration für TypeScript / React-Code (TSX)
% Hinweis: \usepackage{listings} und \usepackage{xcolor} sind in
% praeampel.tex bereits vorhanden und werden hier NICHT erneut geladen.
% =======================================================================

% Farben für die Syntax-Hervorhebung
\definecolor{codebg}{RGB}{246,246,246}
\definecolor{codeframe}{RGB}{200,200,200}
\definecolor{codekeyword}{RGB}{0,0,180}
\definecolor{codecomment}{RGB}{95,140,95}
\definecolor{codestring}{RGB}{170,40,40}

% Eigene Sprache für TypeScript/TSX, da listings sie nicht direkt kennt
\lstdefinelanguage{TypeScript}{
  keywords={
    const, let, var, function, return, if, else, for, while, switch, case,
    break, continue, new, this, typeof, instanceof, in, of, class, extends,
    implements, interface, type, enum, export, import, from, as, default,
    async, await, try, catch, finally, throw, null, undefined, true, false,
    void, never, unknown, any, public, private, protected, readonly, static
  },
  keywordstyle=\color{codekeyword}\bfseries,
  comment=[l]{//},
  morecomment=[s]{/*}{*/},
  commentstyle=\color{codecomment}\itshape,
  string=[b]",
  string=[b]',
  string=[b]`,
  stringstyle=\color{codestring},
  morestring=[b]",
  sensitive=true
}

% Allgemeiner Stil für alle Code-Listings in der Arbeit
\lstdefinestyle{tsxstyle}{
  language=TypeScript,
  backgroundcolor=\color{codebg},
  rulecolor=\color{codeframe},
  frame=single,
  framesep=4pt,
  basicstyle=\ttfamily\footnotesize,
  numbers=left,
  numberstyle=\tiny\color{gray},
  numbersep=8pt,
  breaklines=true,
  breakatwhitespace=false,
  showstringspaces=false,
  tabsize=2,
  captionpos=b,
  xleftmargin=12pt,
  xrightmargin=4pt
}

\lstset{style=tsxstyle}